\documentclass[10pt,a4paper]{amsart}
\usepackage[margin=19mm]{geometry}
\usepackage{amsmath,amssymb,mathtools}
\usepackage[scr=rsfs,cal=euler]{mathalfa}
\usepackage{xfrac}
\usepackage[unicode,hidelinks,bookmarks=false]{hyperref}
\newcommand{\ie}{i.e.\ }
\newcommand{\cf}{cf.\ }
\newcommand{\resp}{respectively\ }
\newcommand{\Subsection}{\subsection}
\providecommand{\nobreakdash}{\nobreak\hbox{-}}
\setlength{\emergencystretch}{2em}
\multlinegap0pt
\usepackage{amssymb}
\usepackage[shortlabels]{enumitem}
\usepackage{float}
\usepackage{tikz}
\newtheorem{lemma}{Lemma}
\title{Generalized Sterboul--Deming Configurations}
\author{Daniel A Jaume, Cristian Panelo, Kevin Pereyra}
\date{Source: arXiv:2609.12128v1 (CC BY 4.0)}
\begin{document}
\maketitle

\noindent\emph{Source and licence.} Generalized Sterboul--Deming Configurations. Version arXiv:2609.12128v1 (CC BY 4.0), distributed by arXiv under the Creative Commons Attribution 4.0 International licence, http://creativecommons.org/licenses/by/4.0/, as recorded in the arXiv licence metadata for this exact version and shown on the abstract page https://arxiv.org/abs/2609.12128v1. Full licence notice: \texttt{licenses/arxiv-2609.12128-CC-BY-4.0.txt}.\par\smallskip

\noindent\textit{Editorial scope.} Counted unit: the article's lemma lem\_ear\_absorption (source0.tex lines 781--799). The article's complete original proof of it is delivered with it (source0.tex lines 801--1045, one \texttt{proof} environment of 245 lines). No mathematics is written, repaired or re-derived: the statement and the proof are the article's own lines, in the article's own notation, and no retained line is substituted or re-flowed.  No further source-local context is needed: every cross-reference of every retained line resolves inside the two delivered spans.  Nothing was omitted from any retained span, so the package carries no omission record.  No retained line contains a citation, so the wrapper carries no bibliography and no bibliography entry.  No retained line contains a figure environment, an image-inclusion command or drawing code, so no image is delivered and the package declares no asset dependency.  Editorial formatting only, in addition: an AMS \texttt{amsart} preamble that restates the article's own declarations and macros used by the retained lines, namely the environments \texttt{lemma} on separate counters, together with the standard packages those lines need.  The retained environments print in this document as Lemma 1 in order of appearance, whereas the article prints the same items with the numbering of its own full document.  The complete native source of the article as distributed in the arXiv e-print for this exact version is bundled unchanged as \texttt{sources/210082/source0.tex}.
\begin{lemma}[Odd-Ear Absorption Lemma]
\label{lem_ear_absorption}
Let \(Q\) be an \(M\)-\(T\)-posy in a graph \(G\), and let
\(R=xRy\) be a simple odd \(M\)-alternating ear relative to \(Q\),
whose internal vertices lie outside \(Q\), and whose first and last
edges have status \(n\). We allow \(x=y\), in which case \(R\) is a
simple closed ear.

Set
\[
X:=V(Q)\cup V(R).
\]
Then there exists a perfect matching \(\widehat M\) of \(G\), agreeing
with \(M\) outside \(G[X]\), and an \(\widehat M\)-\(T\)-posy
\(Q'\subseteq G[X]\) such that
\[
\operatorname{int}R\subseteq V(Q').
\]
\end{lemma}

\begin{proof}
We first isolate the matching surgery used throughout the proof.
Since \(Q\) is an \(M\)-\(T\)-posy, every vertex of \(Q\) is
\(M\)-saturated inside \(Q\). Since \(R\) is an odd alternating ear of
type \(nn\), all internal vertices of \(R\) are paired by \(M\)-edges
of \(R\). Hence every vertex of \(X\) has its \(M\)-mate in \(X\), and
therefore no edge of \(M\) has exactly one endpoint in \(X\).

Consequently, whenever we construct a perfect matching
\(\widehat M_X\) of \(G[X]\), it extends to a perfect matching of \(G\)
by
\[
\widehat M
=
\widehat M_X
\cup
\bigl(M\setminus E(G[X])\bigr).
\]
Thus, in each case below, it is enough to exhibit a \(T\)-posy
\(Q'\subseteq G[X]\) containing \(\operatorname{int}R\), together with
a perfect matching of the residual vertices \(X\setminus V(Q')\).
We give \(Q'\) its canonical perfect matching and match the residual
pieces consecutively along the paths described below.

Suppose first that \(x=y\). Then \(R\) is an \(M\)-blossom with base
\(x\). Since \(Q\) is a \(T\)-posy, a direct inspection of its two
possible geometries yields an \(M\)-blossom \(C\subseteq Q\), with base
\(c\), and a nontrivial simple \(M\)-alternating path \(A\) of type
\(mm\) from \(c\) to \(x\), such that
\[
V(A)\cap V(C)=\{c\}.
\]
Indeed, in the barbell case this follows from the parity of the two
routes from \(x\) to the two bases. In the \(K_4\)-subdivision case,
after relabeling so that two opposite links have type \(mm\), the other
four have type \(nn\), and there is an \(M\)-blossom based at each
branch vertex. The four branch-vertex positions, the two \(mm\)-link
positions, and the four \(nn\)-link positions are all covered by the
same parity routing: for an \(mm\)-link use the odd terminal segment;
for an \(nn\)-link prepend the appropriate adjacent \(mm\)-link to the
even terminal segment. In every case the chosen path meets its blossom
only at the base.

Since \(V(R)\cap V(Q)=\{x\}\), the blossoms \(C\) and \(R\) are
vertex-disjoint and \(A\) meets them only at their bases. Hence
\[
Q':=C\cup A\cup R
\]
is an \(M\)-\(T\)-posy in the barbell case, and no change of matching
is required. Hence assume from now on that \(x\neq y\).

Suppose first that \(Q\) is in the barbell case. Write
\[
Q=C_1\cup P\cup C_2,
\]
where \(C_1,C_2\) are the two blossoms, with bases \(b_1,b_2\), and
\(P=b_1Pb_2\) is the joining path. Up to interchanging the two blossoms
and the two ends of \(R\), there are four possible positions of
\(x\) and \(y\).

If \(x,y\in V(P)\), let \(S=xPy\). If \(\|S\|\) is odd, replace \(S\)
by \(R\). The resulting \(b_1\)-\(b_2\) path remains odd and yields a
new barbell \(T\)-posy. The residual vertices are precisely
\(\operatorname{int}S\), which have even cardinality. If \(\|S\|\) is
even, then
\[
D:=R\cup S
\]
is an odd cycle. Of the two tails \(b_1Px\) and \(yPb_2\), exactly one
has odd length. Joining \(D\) to the blossom at the end of that odd
tail yields a barbell \(T\)-posy. The remaining tail has even length;
together with the unused original blossom it leaves only path pieces
of even order, and these admit consecutive perfect matchings.

If \(x\in V(C_1)\) and \(y\in V(P)\), write
\[
a:=\|b_1Py\|,
\qquad
d:=\|yPb_2\|.
\]
The numbers \(a\) and \(d\) have opposite parity. If \(d\) is even,
then
\[
R\cup yPb_2
\]
is an odd path from \(x\) to \(b_2\), and
\[
C_1\cup R\cup yPb_2\cup C_2
\]
is a barbell with bases \(x\) and \(b_2\). The residual is
\(\operatorname{int}(b_1Py)\), which has even order. If \(d\) is odd,
then \(a\) is even. Let \(E\) be the even \(b_1\)-\(x\) arc of \(C_1\)
and put
\[
D:=R\cup yPb_1\cup E.
\]
Then \(D\) is an odd cycle and
\[
D\cup yPb_2\cup C_2
\]
is a barbell. The complementary \(b_1\)-\(x\) arc of \(C_1\) has odd
length, so its internal vertices have even cardinality. The case
\(x=b_1\) is included by taking \(E\) to be the trivial path.

If \(x,y\in V(C_1)\), let \(E\) and \(F\) be respectively the even and
odd \(x\)-\(y\) arcs of \(C_1\). Then
\[
D:=R\cup E
\]
is an odd cycle. If \(b_1\in V(E)\), then
\[
D\cup P\cup C_2
\]
is a barbell, and the residual is \(\operatorname{int}F\), which has
even order. If \(b_1\in V(F)\), write
\[
F=xFb_1Fy.
\]
Exactly one of the two segments \(xFb_1\) and \(b_1Fy\) has even
length. Choosing the corresponding endpoint \(x\) or \(y\) as the base
of \(D\), and joining it to \(C_2\) through that even segment followed
by \(P\), produces an odd joining path. The unused complementary
segment has odd length and hence an even number of internal vertices.

Finally, if \(x\in V(C_1)\) and \(y\in V(C_2)\), then
\[
Q':=C_1\cup R\cup C_2
\]
is itself a barbell with bases \(x\) and \(y\). Give this barbell its
canonical perfect matching. The only residual vertices are
\(\operatorname{int}P\), and they have even cardinality because
\(\|P\|\) is odd.

This exhausts the barbell case.

Now suppose that \(Q\) is an even subdivision of \(K_4\). Label its
branch vertices \(1,2,3,4\) and its links \(L_{ij}\). All six links
have odd length. Up to an automorphism of \(K_4\), the ends \(x,y\)
lie on the same link, on adjacent links, or on opposite links. Branch
vertices are included in these cases; trivial terminal subpaths of
length \(0\) are allowed.

Assume first that \(x,y\in L_{12}\), in this order, and let
\[
S:=xL_{12}y.
\]
If \(\|S\|\) is odd, replace \(S\) by \(R\). The resulting six links
are again all odd, so the resulting graph is an even subdivision of
\(K_4\). The residual is \(\operatorname{int}S\), which has even
cardinality. If \(\|S\|\) is even, then
\[
D:=R\cup S
\]
is an odd cycle. The two tails \(1L_{12}x\) and \(yL_{12}2\) have
opposite parity. The odd tail joins \(D\) to the appropriate odd cycle
through the remaining four branch vertices, producing a barbell
\(T\)-posy. The complementary tail has even length, and together with
the unused link interiors yields only path components of even order.

Next suppose that \(x\in L_{12}\) and \(y\in L_{13}\). Put
\[
\alpha:=\|1L_{12}x\|\pmod2,
\qquad
\beta:=\|1L_{13}y\|\pmod2.
\]
There are four parity pairs \((\alpha,\beta)\). In the case
\((0,0)\), the odd cycle
\[
R\cup xL_{12}1L_{13}y
\]
is joined through \(L_{14}\) to
\[
L_{23}\cup L_{34}\cup L_{42},
\]
and the residual consists of the interiors of \(xL_{12}2\) and
\(yL_{13}3\). In the other three parity cases the witness is an even
subdivision of \(K_4\): the two relevant odd cycles share one odd link,
the four remaining cycle arms are odd, and the joining path is odd.
For \((1,1)\), \((1,0)\), and \((0,1)\), respectively, the residual is
the interior of \(L_{23}\), \(1L_{12}x\), and \(1L_{13}y\). Thus the
four parity cases are all covered, and every residual component has even
order.

Finally suppose that \(x\in L_{12}\) and \(y\in L_{34}\). Let
\(i,i'\in\{1,2\}\) be chosen so that
\[
\|xL_{12}i\|\equiv1,
\qquad
\|xL_{12}i'\|\equiv0
\pmod2,
\]
and let \(j,j'\in\{3,4\}\) be chosen so that
\[
\|yL_{34}j\|\equiv0,
\qquad
\|yL_{34}j'\|\equiv1
\pmod2.
\]
Set
\[
C=
R\cup xL_{12}i\cup L_{ij}\cup jL_{34}y,
\]
\[
D=
L_{12}\cup L_{1j'}\cup L_{2j'},
\]
and
\[
A=yL_{34}j'.
\]
The cycles \(C\) and \(D\) are odd, they share the odd path
\(xL_{12}i\), and \(A\) is odd. More explicitly, the six effective
links have branch vertices \(x,i,y,j'\) and are
\[
xL_{12}i,
\qquad
yL_{34}j',
\qquad
R,
\qquad
L_{ij}\cup jL_{34}y,
\qquad
L_{ij'},
\qquad
xL_{12}i'\cup L_{i'j'}.
\]
Each has odd length. Hence their union is an even subdivision of
\(K_4\), and therefore a \(T\)-posy containing all internal vertices
of \(R\). The only unused original link is \(L_{i'j}\); its internal
vertices have even cardinality and admit a consecutive perfect
matching.

The same constructions remain valid when an end of \(R\) is a base or
a branch vertex; the corresponding terminal subpath is simply trivial.
Thus all possible positions of the ends of \(R\) are covered.

In every case we have obtained a \(T\)-posy \(Q'\subseteq G[X]\)
containing \(\operatorname{int}R\) and a perfect matching of the
residual vertices. Combining the canonical perfect matching of \(Q'\)
with the residual matching gives a perfect matching
\(\widehat M_X\) of \(G[X]\), and the initial matching surgery extends
it to a perfect matching \(\widehat M\) of \(G\) agreeing with \(M\)
outside \(G[X]\).
\end{proof}

\end{document}
